\documentclass[12pt,aspectratio=1610]{beamer}
\RequirePackage{
    amsmath,
    amssymb,
    calc,
    cancel,
    booktabs,
    color,
    siunitx,
    tikz,
    wrapfig,
    array,
    leftidx,
    float,
    etoolbox,
    fancyhdr,
    longtable,
    hyperref,
    ltcaption,
    ulem,
    wasysym,
    accents,
    listings,
    tabularx,
}

\hypersetup{
    hidelinks,
    breaklinks              = true,
}

\usepackage[final]{pdfpages}
\usepackage[many]{tcolorbox}

\RenewDocumentCommand{\vec}{m}{\mathbf{#1}}

\makeatletter
    \def\new@mathgroup{\alloc@8\mathgroup\mathchardef\@cclvi}
    \patchcmd{\document@select@group}{\sixt@@n}{\@cclvi}{}{}
    \patchcmd{\select@group}{\sixt@@n}{\@cclvi}{}{}
\makeatother

\RequirePackage{mathspec}                                   % includes fontspec
\RequirePackage{polyglossia}                                % multi-language support
\RequirePackage{xunicode}
\setdefaultlanguage{slovak}

% Setup fonts -- see fontspec/mathspec documentation.
\defaultfontfeatures{
    Mapping         = tex-text,
    Scale           = MatchLowercase,
    Ligatures       = TeX
}


\NewDocumentCommand{\labelmath}{m +m}{%
    \begin{equation}%
        #2%
        \label{#1}%
    \end{equation}%
}

\NewDocumentCommand{\labelalign}{m +m}{%
    \begin{align}%
        #2%
        \label{#1}%
    \end{align}%
}

\linespread{1.0}
\setlength{\parindent}{0cm}
\setlength{\parskip}{6pt}
\setlength{\abovedisplayskip}{0mm}
\setlength{\belowdisplayskip}{0mm}
\setlength{\abovedisplayshortskip}{0mm}
\setlength{\belowdisplayshortskip}{0mm}
\setlength{\itemindent}{0pt}
\setlength{\textfloatsep}{0mm}
\setlength{\tabcolsep}{3mm}
\setlength{\LTcapwidth}{0.8\textwidth}
\renewcommand{\arraystretch}{1.2}

\setcounter{secnumdepth}{2}

/home/kvik/dgs/core/tex/math.tex
\DeclareSIUnit\au{AU}
\DeclareSIUnit\pixel{px}
\DeclareSIUnit\lightyear{ly}
\DeclareSIUnit\parsec{pc}
\DeclareSIUnit\earthmass{M_{\earth}}
\DeclareSIUnit\speedoflight{c}
\DeclareSIUnit\foe{foe}
\DeclareSIUnit\year{yr}
\DeclareSIUnit\eur{€}
\DeclareSIUnit\solarmass{M_{\astrosun}}
\DeclareSIUnit\solarluminosity{L_{\astrosun}}
\DeclareSIUnit{\byte}{B}




\linespread{1.0}
\setlength{\parindent}{0cm}
\setlength{\parskip}{6pt}
\setlength{\abovedisplayskip}{0mm}
\setlength{\belowdisplayskip}{0mm}
\setlength{\abovedisplayshortskip}{0mm}
\setlength{\belowdisplayshortskip}{0mm}
\setlength{\itemindent}{0pt}
\setlength{\textfloatsep}{0mm}
\setlength{\tabcolsep}{3mm}
\renewcommand{\arraystretch}{1.2}

\setcounter{secnumdepth}{0}

\NewDocumentCommand{\fspicture}{m O{W} O{black}}{
    {
        \setbeamertemplate{navigation symbols}{}
        \setbeamercolor{background canvas}{bg = #3}
        \begin{frame}[plain]
            \begin{tikzpicture}[remember picture, overlay]
                \node[at=(current page.center)] {
                    \ifstrequal{H}{#2}{                                  
                        \includegraphics[height=\paperheight]{#1}%
                    }{%
                        \includegraphics[width=\paperwidth]{#1}%
                    }
                };
            \end{tikzpicture}
        \end{frame}
    }
}

\NewDocumentCommand{\frejm}{m +m}{
    \begin{frame}
        \frametitle{#1}
        #2
    \end{frame}
}

\NewDocumentCommand{\fragfrejm}{m +m}{
    \begin{frame}[fragile]
        \frametitle{#1}
        #2
    \end{frame}
}

\defbeamertemplate{description item}{align center}{\hfill\insertdescriptionitem\hfill}
\definecolor{desc}{rgb}{0.66, 0, 0}
\definecolor{citem}{rgb}{0.72, 0, 0}
\definecolor{csitem}{rgb}{0.90, 0, 0}
\definecolor{cssitem}{rgb}{1, 0.1, 0.1}
\definecolor{qprimarybg}{rgb}{0.95, 0.95, 0.95}
\definecolor{check}{rgb}{0, 0.8, 0}
\definecolor{coded}{rgb}{0.9, 0.9, 0.9}
\definecolor{todo}{rgb}{1.0, 0.3, 0.3}
\definecolor{model}{rgb}{0.75, 0, 0}

\setbeamertemplate{navigation symbols}{}
\newfontfamily{\semibold}{Segoe UI Semibold}
\RenewDocumentCommand{\emph}{m}{{\semibold#1}}
\NewDocumentCommand{\code}{m}{\textcolor{desc}{\texttt{#1}}}
\NewDocumentCommand{\model}{m}{\colorbox{coded}{\textcolor{model}{\texttt{#1}}}}
\NewDocumentCommand{\todo}{m}{\colorbox{todo}{#1}}

\mode<presentation> {
    \usetheme{Szeged}
    \usecolortheme{beaver}
    
    \usefonttheme{professionalfonts}
    \setallmainfonts{Minion Pro}
    \setmathrm{Minion Pro}
    
    \setsansfont{Segoe UI}
    \setmonofont{Consolas}
    \setbeamercolor*{enumerate item}{fg = citem}
    \setbeamercolor*{enumerate subitem}{fg = csitem}
    \setbeamercolor*{enumerate subsubitem}{fg = cssitem}
    \setbeamercolor*{description item}{fg = desc}
    \setbeamercolor*{itemize item}{fg = citem}
    \setbeamercolor*{itemize subitem}{fg = csitem}
    \setbeamercolor*{itemize subsubitem}{fg = cssitem}
    \setbeamercolor*{palette primary}{fg = red, bg = qprimarybg}
}

\newcommand<>\highlightbox[2]{%
    \alt#3{\makebox[\dimexpr\width-2\fboxsep]{\colorbox{#1}{#2}}}{#2}%
}

\AtBeginSection[]{
    \subsection{\insertsection}
    \begin{frame}
        \vfill
        \centering
        \begin{beamercolorbox}[sep = 18pt, center, shadow = true, rounded = true]{title}
            \usebeamerfont{title}\insertsectionhead%
            \vfill
        \end{beamercolorbox}
        \vfill
    \end{frame}
}

\makeatletter
% Render percent sign with nice font, not ugly Computer modern
    \mathcode`\%="7025

% Fixes mathspec bug -- URL numbers are rendered with wrong font
    \ernewcommand\eu@MathPunctuation@symfont{Latin:m:n}
    \DeclareMathSymbol{,}{\mathpunct}{\eu@MathPunctuation@symfont}{`,}
    \DeclareMathSymbol{?}{\mathpunct}{\eu@MathPunctuation@symfont}{`?}
    \DeclareMathSymbol{.}{\mathord}{\eu@MathPunctuation@symfont}{`.}
    \DeclareMathSymbol{<}{\mathrel}{\eu@MathPunctuation@symfont}{`<}
    \DeclareMathSymbol{>}{\mathrel}{\eu@MathPunctuation@symfont}{`>}
    \DeclareMathSymbol{/}{\mathord}{\eu@MathPunctuation@symfont}{`/}
    \DeclareMathSymbol{;}{\mathpunct}{\eu@MathPunctuation@symfont}{`;}
    \DeclareMathSymbol{(}{\mathopen}{\eu@DigitsArabic@symfont}{`(}
    \DeclareMathSymbol{)}{\mathclose}{\eu@DigitsArabic@symfont}{`)}
    \XeTeXDeclareMathSymbol{^^^^2026}{\mathinner}{\eu@MathPunctuation@symfont}{"2026}[\mathellipsis]
    \DeclareMathSymbol{0}{\mathalpha}{\eu@DigitsArabic@symfont}{`0}
    \DeclareMathSymbol{1}{\mathalpha}{\eu@DigitsArabic@symfont}{`1}
    \DeclareMathSymbol{2}{\mathalpha}{\eu@DigitsArabic@symfont}{`2}
    \DeclareMathSymbol{3}{\mathalpha}{\eu@DigitsArabic@symfont}{`3}
    \DeclareMathSymbol{4}{\mathalpha}{\eu@DigitsArabic@symfont}{`4}
    \DeclareMathSymbol{5}{\mathalpha}{\eu@DigitsArabic@symfont}{`5}
    \DeclareMathSymbol{6}{\mathalpha}{\eu@DigitsArabic@symfont}{`6}
    \DeclareMathSymbol{7}{\mathalpha}{\eu@DigitsArabic@symfont}{`7}
    \DeclareMathSymbol{8}{\mathalpha}{\eu@DigitsArabic@symfont}{`8}
    \DeclareMathSymbol{9}{\mathalpha}{\eu@DigitsArabic@symfont}{`9}
\makeatother


\title{???}
\subtitle{Traja pátrači: záhada koeficientu Γ}
\author{\small \emph{Martin Baláž}}
\institute{MPH seminár}
\date{2021--04--23}

\begin{document}
    {
        \begin{frame}
            \titlepage
        \end{frame}
    }

    \section{Overview}
        \frejm{Odpor vzduchu}{
            Ab initio:
            $$
                F = \Derivative{p}{t} = k \frac{Mv}{t} = 2 e \frac{Mv}{t}
            $$
            \pause
            $$
                M = \rho V = \rho S v t
            $$
            \pause
            $$\vec{F} = c S \rho v^2 \Hat[-1.7mm]{v}$$
            \pause
            $$\vec{F} = -\frac{1}{2} \Gamma S \rho v^2 \Hat[-1.7mm]{v}$$
        }

        \frejm{Koľko je Γ?}{
            Klasické riešenie pre guľu: Γ = \num{0.47}
            \begin{itemize}
                \item ...lenže to platí pre Re = 10000
            \end{itemize}
            \pause
            Čo teraz?
            \begin{itemize}
                \item Karol: CFD
                \item experimentálne korelácie
            \end{itemize}
        }

        \frejm{Drag crisis / Eiffel paradox}{
            \begin{itemize}
                \item Prečo je golfová loptička poďobaná?
                \item Prečo padajúci meteorit v exponenciálne hustnúcej atmosfére zrýchľuje?
            \end{itemize}
        }

        \setbeamersize{description width = 30mm}
        \frejm{Reynoldsovo číslo}{
            $$\mathrm{Re} = \frac{L \rho_\infty v_\infty}{\mu}$$
            kde
            \begin{description}
                \item[$L$] charakteristický rozmer
                \item[$\rho_\infty$] hustota prostredia (unperturbed)
                \item[$v_\infty$] rýchlosť voči prostrediu (unperturbed)
                \item[$\mu$] koeficient dynamickej viskozity
            \end{description}
            \\[5mm]
            \begin{itemize}
                \item inerciálne efekty pre veľké Re
                \item viskózne efekty pre malé Re
            \end{itemize}
        }

        \fspicture{pictures/morrison-drag.png}[][white]

        \frejm{...stále máme problém}{
            \begin{columns}
                \begin{column}{0.7\textwidth}
                    \begin{itemize}
                        \item častica \SI{e6}{\kilo\gram} vo výške \SI{180}{\kilo\metre}
                        \pause
                        \item terminálna rýchlosť \SI{65}{\metre\per\second}?
                    \end{itemize}
                    \pause
                \end{column}
                \begin{column}{0.29\textwidth}
                    \includegraphics[keepaspectratio, width=\textwidth]{pictures/wat.jpg}
                \end{column}
            \end{columns}
            \pause
            Čo s tým?
            \begin{itemize}
                \item model platí pre \emph{nestlačiteľnú} kvapalinu
                \item dva ďalšie efekty
                    \begin{itemize}
                        \item riedka atmosféra vo veľkých výškach -- non-continuum effects
                        \item vysoko nadzvuková rýchlosť a rázové vlny -- 
                    \end{itemize}
            \end{itemize}
        }

        \frejm{Knudsenovo číslo}{
            Pomer strednej voľnej dráhy a charakteristického rozmeru
            $$
                \mathrm{Kn} = \frac{\lambda}{L}
            $$
            Pre ideálny plyn
            $$
                \mathrm{Kn} = \frac{kT}{\sqrt{2} \pi d^2 p L} = \frac{\mu}{\rho L}\sqrt{\frac{\pi m}{2 k T}}
            $$
            \pause
            \begin{description}
                \item[$\mathrm{Kn} < 0.01$]         continuum flow
                \item[$0.01 < \mathrm{Kn} < 0.1$]   slip flow
                \item[$0.1 < \mathrm{Kn} < 10$]     transitional flow
                \item[$10 < \mathrm{Kn}$]           free-molecular flow
            \end{description}
        }

        \frejm{Machovo číslo}{
            Pomer rýchlosti k rýchlosti zvuku v médiu
            \begin{itemize}
                \item a to je pri meteoroidoch dosť problém
                \item model iba s $\mathrm{Re}$ sa nedá použiť
            \end{itemize}

            a platí vzťah...
            $$
                \mathrm{Kn} = \frac{\mathrm{Ma}}{\mathrm{Re}} \sqrt{\frac{\gamma\pi}{2}}
            $$

            \begin{description}
                \item[$\mathrm{Ma} < 0.3$]          incompressible
                \item[$0.3 < \mathrm{Ma} < 1$]      compressible
                \item[$1 < \mathrm{Ma} < 5$]        supersonic
                \item[$5 < \mathrm{Ma}$]            hypersonic
            \end{description}
        }

        \frejm{Singh+2020}{
        }

        \frejm{Implementácia}{
        }

\end{document}
